\documentclass[dvisvgm,tikz,border=3pt]{standalone}
\usepackage{amsmath,amssymb}
\usepackage{pgfplots}
\usepackage{CJKutf8}
\pgfplotsset{compat=1.18}
\usetikzlibrary{arrows.meta,calc,positioning}

\definecolor{themebg}{HTML}{2e362c}
\definecolor{themefg}{HTML}{edf4e8}
\definecolor{themegray}{HTML}{a4af9d}
\definecolor{themecurve}{HTML}{7eb6ff}
\definecolor{themealert}{HTML}{ff7b7b}
\definecolor{themeamber}{HTML}{f5b942}

\pgfdeclarelayer{background}
\pgfdeclarelayer{foreground}
\pgfsetlayers{background,main,foreground}

\begin{document}
\begin{CJK*}{UTF8}{gbsn}
\pagecolor{themebg}
\begin{tikzpicture}[color=themefg, text=themefg]

\begin{axis}[
    width=13.5cm,
    height=8.0cm,
    axis lines=middle,
    axis line style={color=themefg, -{Stealth}},
    tick style={color=themefg},
    tick label style={color=themefg, font=\scriptsize},
    every axis label/.style={color=themefg},
    xmin=0, xmax=16.5,
    ymin=0.5, ymax=3.5,
    xtick={1,3,5,7,9,11,13,15},
    ytick={2},
    yticklabels={$a$},
    clip=false,
    xlabel={$n$},
    ylabel={$x_n$},
    title style={font=\footnotesize\bfseries, at={(0.5,1.06)}, anchor=south},
    title={数列极限 $\varepsilon\text{--}N$ 抓捕带与子列覆盖：$x_n \to a \iff x_{2k} \to a \land x_{2k-1} \to a$},
]

% 1. epsilon 抓捕带 (使用 themecurve!20!themebg 避免白光)
\begin{pgfonlayer}{background}
    \fill[themecurve!20!themebg] (axis cs:0, 1.7) rectangle (axis cs:16, 2.3);
    \draw[themecurve!70!themebg, dashed, line width=1.1pt] (axis cs:0, 2.3) -- (axis cs:16, 2.3) node[right, font=\tiny\bfseries] {$a+\varepsilon$};
    \draw[themecurve!70!themebg, dashed, line width=1.1pt] (axis cs:0, 1.7) -- (axis cs:16, 1.7) node[right, font=\tiny\bfseries] {$a-\varepsilon$};
    \draw[themefg, dashed, line width=1.3pt] (axis cs:0, 2.0) -- (axis cs:16, 2.0) node[right, font=\scriptsize\bfseries] {极限 $a$};
\end{pgfonlayer}

% 2. 门禁阈值 N 虚线
\draw[themealert, line width=1.2pt, dash dot] (axis cs:7.5, 0.6) -- (axis cs:7.5, 3.2);
\node[text=themealert, font=\scriptsize\bfseries, anchor=south] at (axis cs:7.5, 3.25) {门禁阈值 $N$};

% 3. 偶子列 x_{2k} (暖金色菱形)
\addplot[
    only marks,
    mark=diamond*,
    mark size=3.0pt,
    themeamber,
] coordinates {
    (2, 2.8) (4, 2.45) (6, 2.25) (8, 2.15) (10, 2.08) (12, 2.04) (14, 2.02) (16, 2.01)
};

% 4. 奇子列 x_{2k-1} (晶蓝色圆点)
\addplot[
    only marks,
    mark=*,
    mark size=2.6pt,
    themecurve,
] coordinates {
    (1, 1.1) (3, 1.45) (5, 1.68) (7, 1.82) (9, 1.92) (11, 1.96) (13, 1.98) (15, 1.99)
};

% 5. 前景标注
\begin{pgfonlayer}{foreground}
    \node[text=themeamber, font=\scriptsize\bfseries, anchor=west] at (axis cs:2.3, 2.85) {偶子列 $x_{2k}$ 从上方逼近};
    \node[text=themecurve, font=\scriptsize\bfseries, anchor=west] at (axis cs:1.3, 1.1) {奇子列 $x_{2k-1}$ 从下方逼近};

    % 结论框
    \node[text=themefg, font=\scriptsize, anchor=north west] at (axis cs:0.3, 3.4) {
        \begin{tabular}{l}
        \textbf{子列覆盖与同收敛定理}：\\
        1. 原数列收敛 $\implies$ 任意子列必收敛于同一极限 $a$\\
        2. 奇偶两子列收敛于同一极限 $a \implies$ 原数列必收敛于 $a$\\
        3. 若奇偶子列极限不相等 $\implies$ 原数列必发散 (如 $(-1)^n$)
        \end{tabular}
    };
\end{pgfonlayer}

\end{axis}
\end{tikzpicture}
\end{CJK*}
\end{document}
